% Slajd 3 – rozwiązanie jako usługa dla trzech klientów (kryteria: jakość prototypu 20 %, użyteczność 25 %).
\begin{frame}{Rozwiązanie: sprawdzona dostępność jako usługa}
  \begin{columns}[T,onlytextwidth]
    \begin{column}{0.58\textwidth}
      \footnotesize
      \begin{itemize}
        \item \textbf{Mieszkaniec i~turysta (bezpłatnie).} Profil potrzeb zamiast pytania o~niepełnosprawność;
              karta miejsca z~faktami, źródłem, datą i~statusem; mapa barier; asystent w~języku naturalnym,
              który pokazuje tylko fakty z~bazy.
        \item \textbf{Właściciel lokalu (plan Pro).} Panel: potwierdza i~uzupełnia dane, odpowiada na pytania gości,
              widzi, czego szukają i~czego brakuje; odznaka „Potwierdzone przez właściciela”.
        \item \textbf{Miasto i~partnerzy.} Aktualne dane bez utrzymywania bazy, statystyki braków do planowania,
              Open API dla map i~rezerwacji, wdrożenie white-label w~kolejnym mieście.
      \end{itemize}
      \vspace{1mm}
      \begin{exampleblock}{Działa dziś, nie na makietach}
        \footnotesize
        6\,545 miejsc i~27 atrybutów, 11 warstw danych miasta i~OSM, mapa barier, panele właściciela i~administratora,
        4~języki, 503 testy. Aplikacja webowa na telefon i~komputer, bez instalacji.
      \end{exampleblock}
    \end{column}
    \begin{column}{0.4\textwidth}
      \screenshot{mapa-desktop}
      \vspace{0.3mm}
      \muted{Profil „Wózek”: miejsca, warstwy miasta i~OSM, zgłoszone bariery z~wagą i~potwierdzeniami.
        Lista obok mapy jest jej tekstową alternatywą.}
      \vspace{1.5mm}

      \kpi{6\,545}{miejsc}\hfill\kpi{27}{atrybutów}\hfill\kpi{11}{warstw danych}
    \end{column}
  \end{columns}

  \note[item]{Jedno zdanie na klienta: gość dostaje odpowiedź „czy to zadziała dla mnie”, lokal dostaje kanał do gości i~wiedzę, czego im brakuje, miasto dostaje dane, których nie musi utrzymywać.}
  \note[item]{Atrybuty: schody, podjazd, winda, drzwi $\geq$~90~cm, toaleta, przewijak, pętla indukcyjna, odpoczynek, parking OzN, cisza, pies, pokoje dostosowane, recepcja 24~h, języki obsługi.}
  \note[item]{Warstwy: przystanki z~krawężnikiem peronowym, toalety, parkingi OzN, trasy i~infrastruktura rowerowa (dane miasta); windy, krawężniki, przejścia z~oznaczeniami dotykowymi, przechowalnie bagażu, SOR i~apteki 24~h (OSM).}
  \note[item]{Asystent AI jest opcjonalny; w~kontenerze demo działa tryb reguł, więc nic nie zależy od sieci. Udogodnienia hurtowo w~panelu właściciela to zalążek pakietu eventowego.}
\end{frame}
